\section{A matched intervention ladder}
\label{sec:interventions}
We distinguish three choices: where gates create exact zeros, where
training penalizes activation magnitudes, and how the optimizer incorporates
that pressure. Table~\ref{tab:design} names the combinations.
The feed-forward hidden state is $h$. Four-site gating (\Afour) additionally
acts on attention-branch normalization output $a$, feed-forward-branch
normalization output $m$, and the concatenated attention context $z$ before
the output projection. Seven-site gating (\Aseven) adds the post-rotation
queries $q$, post-rotation keys $k$, and values $v$.

\begin{table}[t]
\centering
\caption{Intervention design. The 14M study includes every row. The selected
70M and appendix 410M studies include A0, A1-H, A4+OL1@4, and A7+OL1@7.
For local pressure, dose is $\lambda$; for A4/A7, dose is $\kappa$ with
$\lambda=1$ whenever pressure is present.}
\label{tab:design}
\small
\begin{tabular}{llll}
\toprule
Condition & Gated sites & Pressure sites & Update rule \\
\midrule
A0 & native GELU at $h$ & none & task only \\
A1-H & ReLU at $h$ & none & task only \\
A1-H+L1 & ReLU at $h$ & $h$ & joint-loss L1 \\
A1-H+OL1 & ReLU at $h$ & $h$ & OL1 \\
A4 & $a,m,h,z$ & none & task only \\
A4+OL1@h & $a,m,h,z$ & $h$ & OL1 \\
A4+OL1@4 & $a,m,h,z$ & $a,m,h,z$ & OL1 \\
A7 & $a,m,h,q,k,v,z$ & none & task only \\
A7+OL1@7 & $a,m,h,q,k,v,z$ & all seven & OL1 \\
\bottomrule
\end{tabular}
\end{table}

\paragraph{Gate operators.}
For $\kappa\geq0$, the four common sites use a one-sided threshold and the
additional attention sites use a sign-preserving threshold:
\begin{equation}
  \Gplus(x)=x\,\mathbf{1}[x\geq\kappa],
  \qquad
  \Gpm(x)=x\,\mathbf{1}[|x|\geq\kappa].
  \label{eq:gates}
\end{equation}
At $h$, the gate replaces GELU. Elsewhere it acts on the named tensor.
Gradients pass through retained values. At $\kappa=0$, the added $q,k,v$
gates are identities in both values and gradients, providing a
same-architecture control for the unpressured A4/A7 comparison.
The common one-sided gates remain active at this threshold.

\paragraph{Activation pressure.}
For the tensors captured at the selected layers and sites, the penalty is
$\mathcal{P}=J^{-1}\sum_{j=1}^{J}|A_j|^{-1}\|A_j\|_1$.
Thus each tensor contributes its mean magnitude with equal weight.
The joint-loss baseline applies AdamW to
$\mathcal{L}_{\rm CE}+\lambda\mathcal{P}$. OL1 instead maintains task-only
Adam moments, preconditions the pressure gradient with the task second moment,
removes its component opposing the task update when their dot product is
negative, and caps the weighted pressure-direction norm at the task-direction
norm. Appendix~\ref{app:optimization} specifies the update.

\paragraph{What the contrasts isolate.}
A4 versus A4+OL1@h tests hidden-state pressure under fixed four-site gates.
A4+OL1@h versus A4+OL1@4 expands the pressure sites with the gate set fixed.
A4 versus A7 adds the three attention-operand gates with pressure absent.
Within each of A4 and A7, adding OL1 tests pressure at that configuration's
sites. Comparing these last two effects is descriptive because the pressure
set expands together with the gate set. A7 gates with pressure confined to
$a,m,h,z$ would complete the fixed-pressure comparison.

The A4+OL1@h cohort was originally labeled more broadly in its run metadata.
Artifact inspection established the executed $h$-only pressure set; we use that
operational identity throughout and preserve the original files.
